\subsubsection{P4：动态规划思想、混合整数规划、目标规划与滚动时域重排}

P4 对应动态规划中的“状态更新后继续决策”思想，也把 P2/P3 的静态调度模型推广为带扰动事件的动态取送货时间窗问题。动态事件发生后，系统以事件时刻为滚动边界，将任务划分为固定任务集 $J_{fix}$ 和可重排任务集 $J_{free}$：已经完成、正在执行且不受事件影响的任务保持不变，受临时新增任务、通道封锁、AMR 延误和局部排程调整影响的未来任务进入重优化窗口。这里的关键不是重新从零开始，而是把“事件发生后的 AMR 可用时间、当前位置、剩余任务池”视为新的系统状态，再在这个新状态上继续安排后续决策。这与动态规划中“状态转移后继续求解剩余子问题”的思路是一致的。

\paragraph{混合整数规划建模。}
从运筹学模型看，P4 可抽象为带动态事件约束的混合整数规划。二进制变量 $x_{ki}$ 表示任务 $i$ 是否由 AMR $k$ 执行，$y_{kij}$ 表示同一 AMR 内任务 $i$ 与任务 $j$ 的前后继关系，路径选择变量表示 P1 候选路径库中的具体路径；连续变量 $S_i,C_i$ 表示任务开始和完成时间。模型约束包括任务唯一分配、AMR 序列连续、时间递推、路径可达性、电量安全、通道封锁时间窗避让、AMR 延误窗口互斥，以及 P3 给出的轨迹 footprint 冲突约束。因此 P4 在动态约束下联合决定“谁执行、何时执行、走哪条路径、是否等待”。

\paragraph{目标规划与字典序优化。}
P4 的目标函数采用目标规划和字典序多目标思想。第一层目标是硬可行性，优先使轨迹冲突数、封锁违规数、AMR 延误违规数、缺失路径数和电量违反数为 0；第二层目标是服务水平，最小化新增延期、总延期和高优先级任务延期；第三层目标是运行质量，控制最大完成时间、路径成本、能耗、负载不均衡和被扰动任务数。这样的目标层级避免了用较小运行成本去抵消硬约束违规，也避免了“全局重排虽然目标值低但扰动过大”的不合理方案。

\paragraph{滚动时域启发式求解。}
完整时空 MILP 的变量规模会随 AMR 数、任务数、候选路径数和时间离散粒度快速增长。P4 采用滚动时域启发式：冻结历史计划，构造受影响任务池，通过后悔值插入、冲突感知短名单、有界局部修复和等待调整快速生成可行解。若用课程语言概括，这一过程可以理解为“先根据事件完成一次状态转移，再在新状态上继续选择后续动作”的动态规划思想，加上整数规划建模与目标规划评价。它的核心价值是在实时性、可行性、服务水平和计划稳定性之间取得可解释的折中。
